\input preamble.tex
\Huge Årsprøve i automatiseringssystemer\\ 3AUA Gand VGS 2021

\normalsize

\vskip 10pt 
Kompetansemål:
\begin{itemize}[noitemsep]
\item utføre arbeid på automatiserte anlegg fagmessig, nøyaktig og i overensstemmelse med krav til helse, miljø og sikkerhet og rutiner for kvalitetssikring og internkontroll
\item utføre risikovurdering og vurdere tiltak for ivaretakelse av person– og maskinsikkerhet
\item vurdere hvilke regelverk og normer som gjelder for arbeidet som skal utføres og anvende dette
\item planlegge, utføre, vurdere kvalitet, sluttkontrollere og dokumentere arbeidet
\item planlegge, programmere, montere og idriftsette programmerbare styresystemer
\item endre og tilpasse skjermbilder for grensesnitt mellom menneske og maskin
\item anvende ulike elektroniske kommunikasjonssystemer i automatiserte anlegg
\item tegne, lese og forklare instrumenterte prosessflytskjemaer og bruke annen relevant dokumentasjon for automatiserte anlegg
\item montere, konfigurere, kalibrere og idriftsettelse digitale og analoge målesystemer
\item måle fysiske størrelser i automatiserte anlegg
\item feilsøke og rette feil i automatiserte anlegg
\item bruke gjeldende regelverk og normer for elektriske installasjoner på maskiner
\end{itemize}

Alle ark som leveres inn skal ha elevens navn. \\ 

Oppgave 1-5 skal leveres på papir, etter levering kan eleven ta frem PC og svare på oppgave 6-9.\\

\vskip 2.5pt 

\hrule
\vfil \eject
Oppgave 1 (6p)%Navngi
\vskip 2.5pt 
På P\&ID-en fra forberedelsesdelen brukes flere ulike instrument identifiserings tag. Hva står følgende tag for:\begin{enumerate}
	\item TT \\\\
\begin{tikzpicture}
	\draw[step=0.5cm,gray!20,very thin]  grid (16,1) ;
\end{tikzpicture}
	\item PT\\\\
\begin{tikzpicture}
	\draw[step=0.5cm,gray!20,very thin]  grid (16,1) ;
\end{tikzpicture}
	\item PZH\\\\
\begin{tikzpicture}
	\draw[step=0.5cm,gray!20,very thin]  grid (16,1) ;
\end{tikzpicture}
	\item PZHH\\\\
\begin{tikzpicture}
	\draw[step=0.5cm,gray!20,very thin]  grid (16,1) ;
\end{tikzpicture}
\end{enumerate}


\vskip 2.5pt 


Oppgave 2 (6p) %Definisjoner
\vskip 2.5pt 
a) Hva er nødstopp \\
\vskip 2.5pt 
\begin{tikzpicture}
	\draw[step=0.5cm,gray!20,very thin]  grid (17,4) ;
\end{tikzpicture}
\vskip 2.5pt 
b) Hvor i lovverket stilles det krav til bruk av nødstopp \\
\vskip 2.5pt 
\begin{tikzpicture}
	\draw[step=0.5cm,gray!20,very thin]  grid (17,4) ;
\end{tikzpicture}
\vskip 2.5pt 
\newpage
Oppgave 3 (6p)%Enkel utregninger.
\vskip 2.5pt 
Under den årlige servicen skal det monteres en flowmåler for å kunne måle COP til anlegget. Du har fått i oppdraget å finne k verien for flowmåleren som skal brukes. 
Følgende er et utdrag fra service manual til unisab:
\vskip 15pt 
$\includegraphics[width=6cm]{~/08FlowCalib.png}$\\

\vskip 15pt 
Her ser vi at Unisab kan måle flow ved  å telle pulser. Den har instillingsmuligheter fra 0.01 til 10 liter pr. puls. 
\vskip 5pt 
a) Skisser et måleoppsett der du kan finne k verdien til en turbinflowmåler som skal brues. 

\begin{tikzpicture}
	\draw[step=0.5cm,gray!20,very thin]  grid (17,10) ;
\end{tikzpicture}
\vskip 10pt 
b) Vis hvordan du kan regne ut k-verdien som skal legges inn i unisab III. 

\vskip 2.5pt 
\begin{tikzpicture}
	\draw[step=0.5cm,gray!20,very thin]  grid (17,6) ;
\end{tikzpicture}
\vskip 2.5pt 
\newpage
Oppgave 4 (6p)
\vskip 2.5pt 
Skjemaet i forbindelsesdelen har blant annet følgende symobler. Sett navn ved hvert symbol på hvilken komponent det er symbol for. \\
$\includegraphics[width=6cm]{~/01automatsikring.png}$\\
\vskip 1.5 cm
$\includegraphics[width=6cm]{~/02estop.png}$\\
\vskip 1.5 cm
$\includegraphics[width=6cm]{~/03kabel.png}$\\
\vskip 1.5 cm
$\includegraphics[width=6cm]{~/04magnetventil.png}$\\
\vskip 1.5 cm
$\includegraphics[width=6cm]{~/05Motorverngbryter.png}$\\
\vskip 1.5 cm
$\includegraphics[width=6cm]{~/06PT100.png}$\\
\vskip 1.5 cm
$\includegraphics[width=6cm]{~/07tidsrele.png}$\\
\vskip 1.5 cm
\vskip 2.5pt 
\begin{tikzpicture}
	\draw[step=0.5cm,gray!20,very thin]  grid (17,4) ;
\end{tikzpicture}
\vskip 2.5pt 
\newpage
Oppgave 5 (6p) %Feilsøkings oppgave nødstoppen

\vskip 2.5pt 
Mekanikkeren du jobber med lurer på hva som  er koblet til -B33 (FLOW COP). Tegn og forklar prinsippet en hastighetsbasert flowmåler og en trykkbasert flowmåler. 
\vskip 2.5pt 
\begin{tikzpicture}
	\draw[step=0.5cm,gray!20,very thin]  grid (17,4) ;
\end{tikzpicture}
\vskip 2.5pt 
\vfil \eject

Oppgave 6 (6p) % Tegn og forklar virkemåte til 
\vskip 2.5pt 
Bedriften skal gå over til å standardisere modbus TCP som kommunikasjonssystem. I forbindelse med årligkontroll skal anlegget omkobles. Ta utgangspunkt i prinsippskisse for syringer fra forbindelsesdelen og tegn denne om slik at en bare bruker modbus TCP. 


\vskip 1 cm
Oppgave 7 (6p) % Tegn og forkalr virkemåte. 
\vskip 2.5pt 
Vis hvordan du ville lagt inn en unisab og hentet følgende signaler fra denne inn i codeys
\begin{itemize}[noitemsep]
	\item -B30 Discharge Pressure
	\item -B86 Discharge temperature
	\item Utløpstrykk høy alarm
	\item Utløpstremperatur høy alarm
\end{itemize}


\vskip 1 cm
Oppgave 8 (6p)
\vskip 2.5pt 
Vis hvordan du ville skallert og presentert signalene fra oppgave 7 på en HMI. -B86 skal også logges. 


\vskip 1 cm
Oppgave 9 (12p)

\vskip 5pt 
I forbindelse med oppstart av anlegget starter ikke kompressorene. Mekanikkeren har funnet ut av at feilen ligger i forbindelse med nødstoppreleet. Du får jobben med å finne feilen i forbindelse med nødstoppreleet. 
\vskip 10pt 
Viser hvordan du bille planlagt, gjennomført og dokumentert jobben. 
\end {document}
